\documentclass[../DoAn.tex]{subfiles}
\begin{document}

% Lưu ý: \textbf{Trước khi viết ĐATN, sinh viên cần đọc kỹ hướng dẫn và quy định chi tiết} về cách viết ĐATN trong Phụ lục A. Sinh viên tuân theo mẫu tài liệu này để viết báo cáo đồ án tốt nghiệp, vì tài liệu này đã được căn chỉnh, chỉnh sửa theo đúng chuẩn báo cáo kỹ thuật đồ án tốt nghiệp (ISO 7144:1986). Sinh viên viết trực tiếp vào file này, chỉ chỉnh sửa nội dung, và không viết trên file mới.

% \textbf{Khi đóng quyển ĐATN}, sinh viên cần lưu ý tuân thủ hướng dẫn ở phụ lục A.9

% \textbf{SV cần đặc biệt lưu ý cách hành văn}. Mỗi đoạn văn không được quá dài và cần có ý tứ rõ ràng, bao gồm duy nhất một ý chính và các ý phân tích bổ trợ để làm rõ hơn ý chính. Các câu văn trong đoạn phải đầy đủ chủ ngữ vị ngữ, cùng hướng đến chủ đề chung. Câu sau phải liên kết với câu trước, đoạn sau liên kết với đoạn trước. Trong văn phong khoa học, sinh viên không được dùng từ trong văn nói, không dùng các từ phóng đại, thái quá, các từ thiếu khách quan, thiên về cảm xúc, về quan điểm cá nhân như “tuyệt vời”, “cực hay”, “cực kỳ hữu ích”, v.v. Các câu văn cần được tối ưu hóa, đảm bảo rất khó để thể thêm hoặc bớt đi được dù chỉ một từ. Cách diễn đạt cần ngắn gọn, súc tích, không dài dòng.

% Mẫu ĐATN này được thiết kế phù hợp nhất với đa số các đề tài xây dựng phần mềm ứng dụng. Với các dạng đề tài khác (giải pháp, nghiên cứu, phần mềm đặc thù, v.v.), sinh viên dựa trên cấu trúc và hướng dẫn của báo cáo này để đề xuất và trao đổi với giáo viên hướng dẫn để thiết kế khung báo cáo đồ án cho phù hợp. Sinh viên lưu ý \textbf{trong mọi trường hợp, SV luôn phải sử dụng định dạng báo cáo này, và phải đọc kỹ toàn bộ các hướng dẫn từ đầu tới cuối.} Các hướng dẫn không chỉ áp dụng riêng cho đề tài ứng dụng, mà còn phù hợp với các dạng đề tài khác. Ngoài ra, trong mẫu ĐATN này đã được tích hợp một số hướng dẫn dành riêng cho đề tài nghiên cứu.

% Chương 1 có độ dài từ 3 đến 6 trang với các nội dung sau đây

\section{Đặt vấn đề}
\label{section:1.1}
% Khi đặt vấn đề, sinh viên cần làm nổi bật mức độ cấp thiết, tầm quan trọng và/hoặc quy mô của bài toán của mình.

% Gợi ý cách trình bày cho sinh viên: Xuất phát từ tình hình thực tế gì, dẫn đến vấn đề hoặc bài toán gì. Vấn đề hoặc bài toán đó, nếu được giải quyết, đem lại lợi ích gì, cho những ai, còn có thể được áp dụng vào các lĩnh vực khác nữa không. Sinh viên cần lưu ý phần này chỉ trình bày vấn đề, tuyệt đối không trình bày giải pháp.

Trong bối cảnh đô thị hóa nhanh và mật độ phương tiện gia tăng, bài toán tổ chức và quản lý giao thông ngày càng trở nên phức tạp, đòi hỏi các phương pháp phân tích có căn cứ dữ liệu để hỗ trợ quy hoạch, đánh giá phương án tổ chức giao thông và kiểm chứng các giả thuyết kỹ thuật. Trên thực tế, việc thử nghiệm trực tiếp ngoài hiện trường thường tốn kém, khó kiểm soát biến số, và có thể tiềm ẩn rủi ro về an toàn. Vì vậy, mô phỏng giao thông trở thành một công cụ quan trọng, giúp tái hiện nhiều kịch bản khác nhau trong điều kiện có thể kiểm soát, từ đó hỗ trợ quá trình ra quyết định.

Trong nhóm các công cụ mô phỏng, SUMO (Simulation of Urban MObility) là một nền tảng mô phỏng vi mô được sử dụng rộng rãi, cho phép mô tả mạng đường, luồng phương tiện và cơ chế điều khiển (như tín hiệu đèn) với mức độ chi tiết cao. Kết quả mô phỏng có thể cung cấp dữ liệu theo thời gian về vị trí, tốc độ, hướng di chuyển của phương tiện, cũng như các thông tin liên quan đến trạng thái mạng. Tuy nhiên, dữ liệu đầu ra này thường có khối lượng lớn và giàu tính kỹ thuật, khiến việc quan sát và khai thác hiệu quả trở thành một thách thức, đặc biệt khi người dùng cần đánh giá trực quan cấu trúc không gian hoặc diễn biến ở các khu vực có hình học phức tạp.

Việc quan sát kết quả mô phỏng bằng giao diện 2D truyền thống thường đáp ứng tốt các thao tác cơ bản, nhưng vẫn gặp hạn chế khi cần: (i) nhận thức không gian tại các nút giao nhiều nhánh hoặc có tổ chức làn phức tạp; (ii) theo dõi đồng thời nhiều đối tượng trong bối cảnh mật độ cao; và (iii) trình bày kết quả cho nhóm người dùng không chuyên hoặc trong các hoạt động trao đổi, thuyết minh phương án. Trong các trường hợp này, góc nhìn 3D có thể giúp người quan sát nắm bắt tốt hơn mối quan hệ không gian giữa đường, làn, phương tiện và tín hiệu điều khiển, đồng thời tạo điều kiện cho các hình thức tương tác và trình bày trực quan.

Mặc dù vậy, các giải pháp hiển thị 3D trong thực tế thường gặp những rào cản như phụ thuộc vào bản đồ dựng sẵn (tĩnh), cần nhiều thao tác thủ công khi thay đổi kịch bản, hoặc khó mở rộng khi số lượng đối tượng và loại đối tượng tăng lên. Ngoài ra, việc chuyển dữ liệu mô phỏng sang môi trường 3D phát sinh các vấn đề kỹ thuật như ánh xạ hình học từ mô tả mạng đường (đa tuyến) sang mô hình cảnh, chuyển đổi hệ tọa độ, đồng bộ thời gian giữa bước mô phỏng và khung hình hiển thị, và đảm bảo hiệu năng khi cập nhật số lượng lớn đối tượng theo thời gian.

Do đó, việc xây dựng một hệ thống hiển thị giao thông 3D dựa trên dữ liệu đầu ra của SUMO, có khả năng tự động dựng cảnh theo bản đồ mô phỏng và cập nhật trạng thái theo thời gian, là cần thiết để nâng cao hiệu quả quan sát, phân tích và truyền đạt kết quả mô phỏng. Hệ thống này đóng vai trò cầu nối giữa dữ liệu mô phỏng (mang tính kỹ thuật) và nhu cầu khai thác trực quan, giúp người dùng tiếp cận, kiểm chứng và trình bày kịch bản giao thông một cách trực quan và nhất quán.

\section{Mục tiêu và phạm vi đề tài}
\label{section:1.2}
% Sinh viên trước tiên cần trình bày tổng quan các kết quả của các nghiên cứu hiện nay cho bài toán giới thiệu ở phần \ref{section:1.1} (đối với đề tài nghiên cứu), hoặc về các sản phẩm hiện tại/về nhu cầu của người dùng (đối với đề tài ứng dụng). Tiếp đến, sinh viên tiến hành so sánh và đánh giá tổng quan các sản phẩm/nghiên cứu này.

% Dựa trên các phân tích và đánh giá ở trên, sinh viên khái quát lại các hạn chế hiện tại đang gặp phải. Trên cơ sở đó, sinh viên sẽ hướng tới giải quyết vấn đề cụ thể gì, khắc phục hạn chế gì, phát triển phần mềm \textbf{có các chức năng chính gì}, tạo nên đột phá gì, v.v.

% Trong phần này, sinh viên lưu ý chỉ trình bày tổng quan, không đi vào chi tiết của vấn đề hoặc giải pháp. Nội dung chi tiết sẽ được trình bày trong các chương tiếp theo, đặc biệt là trong Chương 5.

Từ các vấn đề nêu ở phần \ref{section:1.1}, đồ án hướng tới mục tiêu xây dựng một hệ thống hiển thị mô phỏng giao thông trong không gian 3D, trong đó dữ liệu đầu vào được lấy trực tiếp từ các tệp xuất ra của SUMO và/hoặc từ quá trình chạy mô phỏng. Trọng tâm của hệ thống là hình thành một quy trình chuyển đổi dữ liệu mạng đường và dữ liệu mô phỏng sang các đối tượng 3D trong Unity3D một cách nhất quán, từ đó hỗ trợ quan sát diễn biến giao thông theo thời gian trong một môi trường trực quan hơn so với 2D.

Về mục tiêu cụ thể, đồ án tập trung vào ba nhóm mục tiêu chính. Thứ nhất, xây dựng cơ chế nhập và phân tích dữ liệu mô tả mạng đường của SUMO (tệp \texttt{.net.xml}) để phục hồi hình học và cấu trúc topo cơ bản của mạng trong môi trường 3D, bảo đảm các đối tượng được dựng bám sát dữ liệu nguồn. Thứ hai, xây dựng cơ chế thu thập dữ liệu động theo thời gian của mô phỏng (thông qua TraCI hoặc tệp dữ liệu theo thời gian nếu có), nhằm lấy được trạng thái phương tiện và một số thành phần điều khiển (ví dụ tín hiệu đèn) theo từng bước mô phỏng. Thứ ba, thiết kế cơ chế cập nhật và hiển thị trạng thái trong Unity3D theo thời gian, bảo đảm chuyển động thể hiện đúng quỹ đạo tổng quát, đồng thời duy trì khả năng hiển thị ổn định khi số lượng đối tượng tăng.

Trong phạm vi của đề tài, hệ thống tập trung vào các thành phần cốt lõi gồm: (i) dựng cảnh 3D từ dữ liệu mạng đường, bao gồm các đoạn đường và các yếu tố hình học phục vụ quan sát; (ii) biểu diễn phương tiện như các thực thể di chuyển theo trạng thái thời gian; (iii) biểu diễn các tín hiệu điều khiển ở mức phù hợp với dữ liệu có thể truy xuất; và (iv) chuẩn hóa, ánh xạ dữ liệu sang hệ tọa độ của Unity3D để bảo đảm tính nhất quán trong hiển thị. Phần đánh giá được thực hiện thông qua các kịch bản mô phỏng mẫu, nhằm kiểm chứng tính đúng đắn của việc ánh xạ dữ liệu, khả năng tái hiện diễn biến theo thời gian và tính ổn định của hiển thị.

Ngoài phạm vi của đồ án là các bài toán tối ưu điều khiển giao thông, hiệu chỉnh/tham số hóa mô hình mô phỏng của SUMO, hoặc phát triển thuật toán ra quyết định điều khiển phương tiện. Trong phạm vi này, hệ thống đóng vai trò là lớp trực quan hóa và hỗ trợ phân tích dựa trên dữ liệu mô phỏng sẵn có, qua đó góp phần tăng khả năng quan sát và trình bày kịch bản thay vì can thiệp vào thuật toán mô phỏng.

\section{Định hướng giải pháp}
\label{section:1.3}
% Từ việc xác định rõ nhiệm vụ cần giải quyết ở phần \ref{section:1.2}, sinh viên đề xuất định hướng giải pháp của mình theo trình tự sau: (i) Sinh viên trước tiên trình bày sẽ giải quyết vấn đề theo định hướng, phương pháp, thuật toán, kỹ thuật, hay công nghệ nào; Tiếp theo, (ii) sinh viên mô tả ngắn gọn giải pháp của mình là gì (khi đi theo định hướng/phương pháp nêu trên); và sau cùng, (iii) sinh viên trình bày đóng góp chính của đồ án là gì, kết quả đạt được là gì.

% Sinh viên lưu ý không giải thích hoặc phân tích chi tiết công nghệ/thuật toán trong phần này. Sinh viên chỉ cần nêu tên định hướng công nghệ/thuật toán, mô tả ngắn gọn trong một đến hai câu và giải thích nhanh lý do lựa chọn.

Để hiện thực hóa mục tiêu ở phần \ref{section:1.2}, đồ án lựa chọn hướng tiếp cận dựa trên chuỗi công cụ: SUMO dùng để mô phỏng giao thông vi mô và sinh dữ liệu; Python kết hợp TraCI dùng để truy xuất trạng thái mô phỏng theo thời gian; Unity3D dùng để dựng cảnh và hiển thị 3D. Lựa chọn này phù hợp vì SUMO cung cấp dữ liệu mô phỏng có cấu trúc rõ ràng và phổ biến (XML/TraCI), trong khi Unity3D là môi trường phát triển 3D mạnh, hỗ trợ dựng cảnh, quản lý đối tượng, camera và tương tác, qua đó nâng cao khả năng quan sát trực quan.

Trên cơ sở đó, giải pháp được triển khai theo một quy trình tổng quát gồm các thành phần chính. Thứ nhất, mô-đun phân tích mạng đọc tệp \texttt{.net.xml} để trích xuất hình học của các đoạn đường và làn (thường được biểu diễn dưới dạng đa tuyến), cũng như một số quan hệ cấu trúc cần thiết phục vụ dựng cảnh. Kết quả được chuyển thành một biểu diễn trung gian phù hợp để sinh các đối tượng trong Unity3D. Thứ hai, mô-đun thu thập dữ liệu động kết nối với phiên mô phỏng SUMO thông qua TraCI để lấy trạng thái phương tiện theo từng bước mô phỏng, bao gồm vị trí, hướng chuyển động và các thuộc tính trạng thái cơ bản. Thứ ba, mô-đun ánh xạ và đồng bộ chuyển đổi hệ tọa độ từ không gian mô phỏng sang không gian Unity3D, đồng thời tổ chức cơ chế đồng bộ thời gian giữa bước mô phỏng và chu kỳ cập nhật khung hình nhằm giảm sai lệch cảm nhận khi hiển thị.

Trong quá trình triển khai, một yêu cầu quan trọng là bảo đảm tính nhất quán của dữ liệu giữa các mô-đun: dữ liệu hình học mạng đường cần dùng chung chuẩn tọa độ với dữ liệu phương tiện; các đối tượng được tạo ra trong Unity3D cần có vòng đời phù hợp với trạng thái mô phỏng (tạo mới khi phương tiện xuất hiện, cập nhật khi di chuyển và loại bỏ khi rời mạng). Đồng thời, để đảm bảo khả năng mở rộng, giải pháp ưu tiên tổ chức dữ liệu theo hướng tách biệt phần đọc/chuẩn hóa dữ liệu với phần hiển thị, giúp giảm phụ thuộc và thuận tiện khi thay đổi kịch bản mô phỏng.

Đóng góp chính của đồ án là đề xuất và hiện thực một quy trình chuyển đổi dữ liệu mô phỏng từ SUMO sang môi trường hiển thị 3D trong Unity3D theo hướng tự động, giảm thao tác thủ công khi thay đổi kịch bản. Bên cạnh đó, đồ án xây dựng được cơ chế dựng cảnh từ dữ liệu mạng đường và cơ chế cập nhật trạng thái phương tiện theo thời gian để tái hiện diễn biến giao thông trong không gian 3D. Kết quả thử nghiệm trên các kịch bản mô phỏng mẫu cho thấy hệ thống có thể dựng được mạng đường từ dữ liệu đầu vào và biểu diễn chuyển động phương tiện theo dữ liệu thu từ SUMO một cách ổn định, tạo cơ sở cho việc quan sát và phân tích trực quan.
\section{Bố cục đồ án}
\label{section:1.4}
% Phần còn lại của báo cáo đồ án tốt nghiệp này được tổ chức như sau. 

% Chương 2 trình bày về v.v. 

% Trong Chương 3, em/tôi giới thiệu về v.v.

% \textbf{Chú ý:} Sinh viên cần viết mô tả thành đoạn văn đầy đủ về nội dung chương. Tuyệt đối không viết ý hay gạch đầu dòng. Chương 1 không cần mô tả trong phần này. 

% Ví dụ tham khảo mô tả chương trong phần bố cục đồ án tốt nghiệp: Chương *** trình bày đóng góp chính của đồ án, đó là một nền tảng ABC cho phép khai phá và tích hợp nhiều nguồn dữ liệu, trong đó mỗi nguồn dữ liệu lại có định dạng đặc thù riêng. Nền tảng ABC được phát triển dựa trên khái niệm DEF, là các module ngữ nghĩa trợ giúp người dùng tìm kiếm, tích hợp và hiển thị trực quan dữ liệu theo mô hình cộng tác và mô hình phân tán.

% \textbf{Chú ý:} Trong phần nội dung chính, mỗi chương của đồ án nên có phần Tổng quan và Kết chương. Hai phần này đều có định dạng văn bản “Normal”, sinh viên không cần tạo định dạng riêng, ví dụ như không in đậm/in nghiêng, không đóng khung, v.v. 

% Trong phần Tổng quan của chương N, sinh viên nên có sự liên kết với chương N-1 rồi trình bày sơ qua lý do có mặt của chương N và sự cần thiết của chương này trong đồ án. Sau đó giới thiệu những vấn đề sẽ trình bày trong chương này là gì, trong các đề mục lớn nào.

% Ví dụ về phần Tổng quan: Chương 3 đã thảo luận về nguồn gốc ra đời, cơ sở lý thuyết và các nhiệm vụ chính của bài toán tích hợp dữ liệu. Chương 4 này sẽ trình bày chi tiết các công cụ tích hợp dữ liệu theo hướng tiếp cận “mashup”. Với mục đích và phạm vi của đề tài, sáu nhóm công cụ tích hợp dữ liệu chính được trình bày bao gồm: (i) nhóm công cụ ABC trong phần 4.1, (ii) nhóm công cụ DEF trong phần 4.2, nhóm công cụ GHK trong phần 4.3, v.v.

% Trong phần Kết chương, sinh viên đưa ra một số kết luận quan trọng của chương. Những vấn đề mở ra trong Tổng quan cần được tóm tắt lại nội dung và cách giải quyết/thực hiện như thế nào. Sinh viên lưu ý không viết Kết chương giống hệt Tổng quan. Sau khi đọc phần Kết chương, người đọc sẽ nắm được sơ bộ nội dung và giải pháp cho các vấn đề đã trình bày trong chương. Trong Kết chương, Sinh viên nên có thêm câu liên kết tới chương tiếp theo.

% Ví dụ về phần Kết chương: Chương này đã phân tích chi tiết sáu nhóm công cụ tích hợp dữ liệu. Nhóm công cụ ABC và DEF thích hợp với những bài toán tích hợp dữ liệu phạm vi nhỏ. Trong khi đó, nhóm công cụ GHK lại chứng tỏ thế mạnh của mình với những bài toán cần độ chính xác cao, v.v. Từ kết quả nghiên cứu và phân tích về sáu nhóm công cụ tích hợp dữ liệu này, tôi đã thực hiện phát triển phần mềm tự động bóc tách và tích hợp dữ liệu sử dụng nhóm công cụ GHK. Phần này được trình bày trong chương tiếp theo – Chương 5.

Phần còn lại của báo cáo đồ án tốt nghiệp này được tổ chức như sau: Chương 2 trình bày tổng quan các công cụ và nền tảng được sử dụng trong hệ thống, bao gồm SUMO, TraCI và Unity3D, đồng thời nêu các đặc trưng dữ liệu đầu vào/đầu ra liên quan đến quá trình mô phỏng và trực quan hóa. Chương 3 trình bày phân tích và thiết kế hệ thống, tập trung vào kiến trúc tổng thể, quy trình chuyển đổi dữ liệu và các mô-đun chính phục vụ dựng cảnh và cập nhật trạng thái theo thời gian. Chương 4 trình bày kết quả thực nghiệm, mô tả các kịch bản mô phỏng mẫu, kết quả dựng cảnh và quan sát, đồng thời thảo luận một số đánh giá về tính đúng đắn và tính ổn định của hiển thị. Cuối cùng, Chương 5 tổng kết các kết quả đạt được, các đóng góp chính của đồ án và đề xuất hướng phát triển trong tương lai.

\end{document}